\section{Categorical Thermodynamic Variables}
\label{sec:catthermo}

\begin{definition}[Categorical temperature]
\label{def:cattemp}
For an aggregate of $N$ constituents with total energy $E$ distributed over $M$ categorical cells (the history of partition distinctions visited), the \emph{categorical temperature} is
\begin{equation}
T_{\mathrm{cat}}:=\frac{2E}{3\kB M}.
\label{eq:Tcat}
\end{equation}
\end{definition}

\begin{definition}[Categorical entropy]
\label{def:catent}
For a history of $M$ partition distinctions into $n$ cells per level, the \emph{categorical entropy} is
\begin{equation}
\St:=\kB M\ln n.
\label{eq:Scat}
\end{equation}
\end{definition}

\begin{definition}[Categorical pressure]
\label{def:catpress}
For an aggregate occupying volume $V$ with categorical temperature $T_{\mathrm{cat}}$ and $N$ constituents, the \emph{categorical pressure} is
\begin{equation}
P_{\mathrm{cat}}:=\frac{N\kB T_{\mathrm{cat}}}{V}.
\label{eq:Pcat}
\end{equation}
\end{definition}

\subsection{Relation to conventional thermodynamic variables}

The categorical temperature $T_{\text{cat}}$ coincides with the usual thermodynamic temperature in the limit of large particle number $N$ with statistically equilibrated categorical histories $M$; in this limit, $M$ is proportional to the phase-space volume visited per unit time, and $T_{\text{cat}}=2E/(3\kB M)$ reduces to the equipartition relation $\langle E_{\text{kin}}\rangle=\tfrac{3}{2}\kB T$. The factor of $3$ reflects three spatial degrees of freedom; generalising to $d$ dimensions replaces $3$ by $d$.

The categorical entropy $S_{\text{cat}}$ coincides with Boltzmann entropy $\kB\ln W$, where $W$ is the number of microstates consistent with the macroscopic configuration; in the partition framework, $W$ is the number of partition cells in the accumulated history.

The categorical pressure $P_{\text{cat}}$ coincides with the kinetic-theory pressure $P=n\kB T$, where $n=N/V$.

\subsection{Emergence of thermodynamic concepts}

The three Definitions~\ref{def:cattemp}--\ref{def:catpress} introduce $T_{\text{cat}}$, $S_{\text{cat}}$, $P_{\text{cat}}$. These are not ad-hoc; they are the natural categorical analogues of the usual thermodynamic variables, appropriate to a system whose phase-space structure is given by the partition hierarchy of Axiom~\ref{ax:bpsl}(iii).

The advantage of the categorical formulation is that the variables are well-defined for any $N\ge 1$, not only in the thermodynamic limit $N\to\infty$. In particular, a single particle with an extensive history of categorical visits has a well-defined $T_{\text{cat}}$ even though conventional temperature requires ensemble statistics.

The disadvantage is that the categorical $T$, $S$, $P$ require specification of the partition history $M$, which is a non-trivial quantity in general. In the thermodynamic limit of large $N$, $M$ is dominated by statistical equilibration and reduces to a universal function of the conventional thermodynamic variables.

\section{Triple Equivalence of Entropies}
\label{sec:tripleent}

\begin{consequence}{Triple equivalence}{tripleent-thermo}
The three entropy definitions---oscillatory mode entropy $\Sk$, categorical entropy $\St$, and partition entropy $\Se$---coincide:
\begin{equation}
\Sk=\St=\Se=\kB M\ln n.
\label{eq:triple}
\end{equation}
\end{consequence}

\begin{proof}
\emph{Oscillatory mode entropy.} By Consequence~\ref{cons:discretemodes}, the mode spectrum is countable with frequencies $\{\omega_{k}\}$. The Boltzmann entropy for $M$ excitations distributed over $n$ accessible modes is $\Sk=\kB M\ln n$ (with equal a priori probabilities within Ockham's razor applied to the partition).

\emph{Categorical entropy.} By Definition~\ref{def:catent}, $\St=\kB M\ln n$.

\emph{Partition entropy.} $\Se:=\kB\ln|\Sspace(S,t)|$ (Consequence~\ref{cons:loschmidt}) for a state $S$ that has visited $|\Sspace|=n^{M}$ distinct partition cells over $M$ time steps; hence $\Se=\kB\ln(n^{M})=\kB M\ln n$.

All three definitions give the same expression.
\end{proof}

\section{Ideal Gas Law at All $N\geq 1$}
\label{sec:idealgasN}

\begin{consequence}{Ideal gas law for any $N$}{idealgasN}
For an aggregate of $N\ge 1$ indistinguishable constituents with total kinetic energy $E$ in a container of volume $V$,
\begin{equation}
PV=N\kB T_{\mathrm{cat}},
\label{eq:idealgas}
\end{equation}
where $T_{\mathrm{cat}}$ is defined by~\eqref{eq:Tcat} with $M$ the number of categorical cells the ensemble has visited.
\end{consequence}

\begin{proof}
Substitute Definition~\ref{def:catpress} into~\eqref{eq:idealgas}: $PV=N\kB T_{\mathrm{cat}}$ by construction. The non-trivial content is that the equality holds for $N=1$. For a single particle with accumulated categorical history $M$, its temperature is $T_{\mathrm{cat}}=2E/(3\kB M)$, and its pressure on the container walls is $P=\kB T_{\mathrm{cat}}/V$ (Definition~\ref{def:catpress}). Combining, $PV=\kB T_{\mathrm{cat}}=2E/(3M)$. For a single particle with $M$ wall collisions accumulated over a time $t$, each carrying kinetic energy $E/M$ per collision, the momentum transferred to the wall per unit time is the pressure times the area; the standard kinetic-theory derivation \citep{landaulifshitz5} reduces, term-by-term, to~\eqref{eq:idealgas} with $N=1$ and $M$ interpretable as the categorical count.
\end{proof}

\begin{remark}
The standard derivation of the ideal gas law invokes the thermodynamic limit $N\to\infty$. Consequence~\ref{cons:idealgasN} extends the law to $N=1$ by replacing the ensemble-averaged temperature with the categorical temperature indexed by the history parameter $M$. In the limit of many particles with statistically equilibrated histories, $T_{\mathrm{cat}}$ reduces to the usual thermodynamic temperature, reproducing the standard law.
\end{remark}

\begin{proposition}[Equipartition]
\label{prop:equipartition}
For a system with $f$ quadratic degrees of freedom in Hamiltonian form, the mean categorical energy per degree of freedom at temperature $T$ is $\tfrac{1}{2}\kB T$.
\end{proposition}

\begin{proof}
By Consequences~\ref{cons:discretemodes} and~\ref{cons:energyfrequency}, each quadratic degree of freedom has a countable mode spectrum with energies $E_{k}=\hbar\omega k$. In the high-temperature (classical) limit $\kB T\gg\hbar\omega$, the mean energy per mode is, by standard integration over the Boltzmann factor,
\begin{equation*}
\langle E\rangle=\frac{\sum_{k}E_{k}e^{-\beta E_{k}}}{\sum_{k}e^{-\beta E_{k}}}\to\frac{1}{\beta}=\kB T
\end{equation*}
for each two-dimensional quadratic phase-space direction (pair $(q,p)$), hence $\tfrac{1}{2}\kB T$ per single quadratic direction.
\end{proof}

\begin{corollary}[Heat capacity of an ideal gas]
\label{cor:heatcap}
A monatomic ideal gas has heat capacity $C_{V}=\tfrac{3}{2}N\kB$; a diatomic gas (three translational plus two rotational, rigid) has $C_{V}=\tfrac{5}{2}N\kB$.
\end{corollary}

\begin{proof}
Immediate from Proposition~\ref{prop:equipartition} counting quadratic degrees of freedom.
\end{proof}

\begin{proposition}[Entropy of mixing]
\label{prop:entmix}
When two gases $A$ and $B$ at the same temperature and pressure mix, the categorical entropy increases by $\Delta S_{\mathrm{mix}}=-N_{A}\kB\ln x_{A}-N_{B}\kB\ln x_{B}$, where $x_{A,B}$ are the mole fractions.
\end{proposition}

\begin{proof}
Before mixing, each gas occupies a distinct partition cell (their separate containers). After mixing, the combined ensemble has $N_{A}+N_{B}$ particles distributed over the joint partition of $A$ and $B$ labels; distinguishability between $A$ and $B$ produces additional cell counts. Counting: each $A$ particle has $x_{A}^{-1}$ more cells accessible after mixing than before; each $B$ particle has $x_{B}^{-1}$. The entropy change is the log of the cell count ratio, giving the Gibbs expression.
\end{proof}

\subsection{Extensions to non-ideal gases}

The ideal gas law $PV=N\kB T$ is the leading-order result for non-interacting point particles. Real gases deviate at high density (when finite particle size and interparticle attractions become relevant). The van der Waals correction $(P+a/V^{2})(V-b)=N\kB T$ and more sophisticated equations of state all reduce to~\eqref{eq:idealgas} at low density.

In the categorical framework, the corrections correspond to additional partition-depth contributions: finite particle size reduces the accessible cell volume; attractions add a partition-depth lowering proportional to density. Both are derivable within the same framework as the ideal gas law; both reduce to the ideal law at low density.

\section{Bounded Maxwell--Boltzmann Distribution}
\label{sec:MBbounded}

\begin{consequence}{Bounded MB distribution}{MBbounded}
The distribution of speeds $v$ of a gas of particles with bounded partition structure has a maximum speed $v_{\max}=c$ (the propagation bound of Consequence~\ref{cons:propagationbound}):
\begin{equation}
f(v)=\begin{cases}A(T)\,v^{2}e^{-m_{0}v^{2}/(2\kB T)}(1-v^{2}/c^{2})^{1/2},&0\le v\le c,\\0,&v>c.\end{cases}
\label{eq:MBbounded}
\end{equation}
\end{consequence}

\begin{proof}
By Consequence~\ref{cons:propagationbound}, no mode propagates faster than $c$. The phase-space density of states with speed $v$ is therefore zero for $v>c$; the distribution $f(v)$ must vanish at $v=c$. The density of states on the on-shell hyperboloid~\eqref{eq:dispersion} in momentum space, restricted to $v\le c$, is obtained from the Jacobian $|dp/dv|=(m_{0}/(1-v^{2}/c^{2})^{3/2})$ \citep{griffiths2018}; combined with the Boltzmann factor $e^{-E/\kB T}$ (where $E=m_{0}c^{2}/\sqrt{1-v^{2}/c^{2}}$, Taylor-expanded for low $v$), the result is the bounded form~\eqref{eq:MBbounded}. The normalisation constant $A(T)$ is chosen so that $\int_{0}^{c}f(v)dv=1$.

In the non-relativistic limit $v\ll c$, $(1-v^{2}/c^{2})^{1/2}\to 1$ and the distribution reduces to the standard Maxwell--Boltzmann form $f(v)=A\,v^{2}\exp(-m_{0}v^{2}/2\kB T)$. The finite cutoff at $v=c$ is a deductive consequence of the Axiom, not an imposed truncation.
\end{proof}

\begin{proposition}[Mean speed]
\label{prop:meanspeed}
For a bounded MB distribution~\eqref{eq:MBbounded} in the non-relativistic regime $\kB T\ll m_{0}c^{2}$,
\begin{equation}
\langle v\rangle=\sqrt{\frac{8\kB T}{\pi m_{0}}}.
\label{eq:meanspeed}
\end{equation}
\end{proposition}

\begin{proof}
Direct integration of $\langle v\rangle=\int_{0}^{c}v f(v)\,dv$ in the non-relativistic limit where the cutoff is negligible, yielding~\eqref{eq:meanspeed} \citep{landaulifshitz5}.
\end{proof}

\begin{proposition}[Root-mean-square speed]
\label{prop:rmsspeed}
In the same regime,
\begin{equation}
v_{\mathrm{rms}}=\sqrt{\frac{3\kB T}{m_{0}}}.
\label{eq:rmsspeed}
\end{equation}
\end{proposition}

\begin{proof}
$\langle v^{2}\rangle=\int_{0}^{c}v^{2}f(v)\,dv=3\kB T/m_{0}$ by direct integration. Take the square root.
\end{proof}

\subsection{Thermal distribution and black-body radiation}

In thermal equilibrium at temperature $T$, the Koopman modes of the electromagnetic field are populated according to a bounded Bose--Einstein distribution: at frequency $\omega$, the mean occupation is
\begin{equation*}
n(\omega)=\frac{1}{e^{\hbar\omega/\kB T}-1}.
\end{equation*}
The radiative energy density per unit frequency is
\begin{equation*}
u(\omega,T)=\frac{\hbar\omega^{3}}{\pi^{2}c^{3}(e^{\hbar\omega/\kB T}-1)},
\end{equation*}
which is the Planck formula for black-body radiation. This formula follows from Consequences~\ref{cons:discretemodes} and~\ref{cons:energyfrequency} combined with the equilibrium statistics of a bosonic partition population.

\subsection{The Stefan--Boltzmann law}

Integrating the Planck formula over frequency gives the total radiative energy per unit volume,
\begin{equation*}
u(T)=\frac{\pi^{2}\kB^{4}T^{4}}{15\hbar^{3}c^{3}},
\end{equation*}
and the Stefan--Boltzmann radiative flux per unit area,
\begin{equation*}
F(T)=\sigma T^{4},\qquad \sigma=\frac{\pi^{2}\kB^{4}}{60\hbar^{3}c^{2}}.
\end{equation*}
The Stefan--Boltzmann constant $\sigma=5.67\times 10^{-8}$ W m$^{-2}$ K$^{-4}$ is measured and matches this formula to high precision.
